\documentclass[conference]{IEEEtran}
\IEEEoverridecommandlockouts

\usepackage{cite}
\usepackage{amsmath,amssymb,amsfonts}
\usepackage{algorithm}
\usepackage{algorithmic}
\usepackage{graphicx}
\usepackage{textcomp}
\usepackage{xcolor}
\usepackage{booktabs}
\usepackage{multirow}
\usepackage{array}
\usepackage{url}
\usepackage{hyperref}

\def\BibTeX{{\rm B\kern-.05em{\sc i\kern-.025em b}\kern-.08em
    T\kern-.1667em\lower.7ex\hbox{E}\kern-.125emX}}

\begin{document}

\title{Toward Learning Progression-Grounded Code Comprehension Scaffolding:
A Research Framework Connecting Science Education Assessment,
Automated Scoring, and LLM Evaluation}



\maketitle

%-------------------------------------------------------------------
\begin{abstract}
The deployment of large language model (LLM) tools in computing education has
substantially outpaced the psychometric and pedagogical frameworks needed to
evaluate them. Chang et al.~\cite{chang2024survey} reviewed 46~LLM benchmarks
and confirmed that none assessed developmental appropriateness or equity. This
paper presents a midterm research report on a framework that bridges three
converging research bodies: Learning Progression (LP) theory~\cite{jin2019validation},
automated LP-based scoring~\cite{jin2025automated}, and a comprehensive LLM
evaluation survey~\cite{chang2024survey}. The system, grounded in the
Cognitive-Preference Adaptive Learning (CPAL) architecture, applies LP theory to
code comprehension via a nine-component multi-modal pipeline integrating CodeBERT
embeddings, a Hyperspherical Variational Self-Attention Autoencoder (HVSAE), COKE
cognitive-state tracking, Nestor Bayesian personality profiling, a four-level
hierarchical RL controller, and three dynamically updated knowledge graphs. The
pipeline exposes a FastAPI REST interface, orchestrates nine parallel per-turn
computations in 425--650\,ms, and implements ten adaptive personalization features.
Preliminary results across 10~student conversations (45~interaction turns)
demonstrate 100\% error detection accuracy, a 42.8\% improvement in evaluation
accuracy using HVSAE semantic encoding over baseline keyword matching, and a mean
concentration parameter $\kappa=91.1\%$ indicating highly discriminative latent
representations. Four explicit Research Questions drawn from LP theory structure
a twelve-week remaining programme culminating in a five-step LLM comparison
pipeline and QWK-based ranking system.
\end{abstract}

\begin{IEEEkeywords}
Learning Progressions, Code Comprehension, Large Language Models, Adaptive
Scaffolding, HVSAE, COKE, Automated Scoring, Quadratic Weighted Kappa,
Educational Data Mining, Equity
\end{IEEEkeywords}

%===================================================================
\section{Introduction}
\label{sec:intro}

The rapid adoption of AI coding assistants in computing classrooms has produced a
fundamental mismatch between deployment velocity and evaluative rigor. Students
write working code but cannot explain its internal reasoning mechanisms. LLM tools
enter classrooms with no pedagogical validation, and no principled theory yet
exists to evaluate whether AI-generated explanations develop reasoning ability
rather than merely supplying syntactically correct answers~\cite{chang2023survey}.

Three empirical facts sharpen this diagnosis. First, Chang et al.~\cite{chang2023survey}
catalogued 46~LLM benchmarks and confirmed that none assessed developmental
appropriateness. Second, Jin et al.~\cite{jin2025automated} showed that 25--32\%
of correctly-reasoning student responses are misclassified due to informal language,
and that one assessment item produced a QWK of~0.00 without item context, rising
to~0.68 with context (Model~B strategy). Third, the SOLO taxonomy~\cite{biggs1982solo},
applied in CS education since 2006~\cite{lister2006solo}, has never been subjected
to Partial Credit Model (PCM) threshold-ordering verification.

Learning Progression (LP) theory provides the missing construct infrastructure.
Jin et al.~\cite{jin2019validation} established a five-stage validation
framework grounded in Kane's~\cite{kane2006validation} validity argument structure
and validated through PCM threshold ordering.

This paper makes four contributions: (1)~adapts LP theory from science education
to code comprehension; (2)~presents a working multi-modal personalized-learning
pipeline; (3)~reports preliminary experimental results across 45~interaction turns;
and (4)~introduces a twelve-week remaining-work programme for the first
psychometrically rigorous evaluation of LLM-generated educational
explanations~\cite{patel2025llmeval}.

%===================================================================
\section{Research Questions}
\label{sec:rqs}

\textbf{RQ1} --- \textit{Can a PCM-validated LP for code comprehension be
constructed?} Specifically: can a four-level LP for recursive reasoning be
empirically validated via PCM threshold ordering~\cite{jin2019validation} such that
all inter-level thresholds are correctly ordered?

\textbf{RQ2} --- \textit{How does informal language affect LP-level classification
of code explanations by LLMs?} Specifically: what proportion of correctly-reasoning
student code explanations are misclassified due to informal expression (e.g.,
``the function remembers where it was'' for stack-frame preservation)? This extends
the 25--32\% finding of Jin et al.~\cite{jin2025automated} to code comprehension.

\textbf{RQ3} --- \textit{Does LP-grounded item description (Model~B context)
improve developmental appropriateness of LLM explanations?} Specifically: does
prepending a structured prompt encoding the target LP level,
knowledge-component prerequisites, and the student's COKE cognitive state produce
explanations scoring higher on the four-level LP rubric than question-only (Model~A)
prompts?

\textbf{RQ4} --- \textit{Do LP-calibrated explanations equitably serve diverse
learners?} Specifically: does Differential Item Functioning (DIF) analysis across
subgroups---native English speakers vs.\ emergent bilinguals; novice vs.\
experienced programmers---reveal systematic LP-level underscoring for any subgroup?

%===================================================================
\section{Literature Review}
\label{sec:lit}

\subsection{Learning Progressions: Theory and Psychometric Validation}

Learning progressions (LPs) are empirically validated sequences of reasoning levels
describing student movement from naive to expert
understanding~\cite{jin2019validation}. Jin et al.~\cite{jin2019validation}
established the gold standard for LP validation through a five-stage framework:
Development, Scoring, Generalisation, Extrapolation, and Use. The decisive
statistical test is PCM threshold ordering: disordered thresholds require LP
revision before any downstream use. The SOLO taxonomy~\cite{biggs1982solo},
applied in CS education since Lister et al.~\cite{lister2006solo}, has never
undergone PCM verification, leaving SOLO-based instructional inference on an
unconfirmed psychometric foundation.

\subsection{Automated Scoring and LP-Based Assessment}

Jin et al.~\cite{jin2025automated} investigated transformer-based automated scoring
for LP-grounded constructed-response items. Model~A achieved QWK~$=0.46$, $0.50$,
and $0.00$ on three of four items, failing the operational threshold of
0.70~\cite{williamson2012framework}. Model~B raised these to~$0.89$, $0.83$,
and $0.68$. Three failure modes inform this research: informal language (25--32\%
misclassification), terminological ambiguity, and systematic bilingual underscoring.

\subsection{LLM Evaluation: Scope and Structural Gaps}

Chang et al.~\cite{chang2023survey} surveyed 46~LLM benchmarks and found that none
assessed developmental appropriateness or equity. LLMs showed limited ability in
semantic similarity and paraphrase tasks---precisely the skills required to
recognise informal-but-correct reasoning.

\subsection{Knowledge Graphs and Ontologies in CS Education}

Dessi et al.~\cite{dessi2025csekg} introduced CSE-KG~2.0, a large-scale knowledge
graph from 14.5~million CS papers with prerequisite relationships and concept
taxonomy. Abu-Salih and Alotaibi~\cite{abusalih2024kg} reviewed five KG application
domains in education. The CPAL framework integrates CSE-KG~2.0 and augments it
with dynamically updated student-specific and pedagogical KGs following the
bidirectional update mechanism of Shengju et al.~\cite{shengju2025cross}.

\subsection{Knowledge Tracing and Cognitive Diagnosis}

Sun et al.~\cite{sun2024tcl4kt} proposed TCL4KT, a weighted heterogeneous
graph-based three-view contrastive learning framework for knowledge tracing.
Wei et al.~\cite{wei2025dakt} introduced DAKT-SSO-MRC, integrating educational
psychology with deep learning. The DINA model~\cite{delatorre2009dina} provides
interpretable mastery estimation via slip and guess probabilities estimated by
EM algorithm on ASSISTments data. Wan et al.~\cite{wan2023rl} demonstrated RL
agents trained on knowledge tracing simulators significantly improve student
knowledge states.

\subsection{Psychological Profiling and Behavioral Analysis}

Nadimpalli et al.~\cite{nadimpalli2025nestor} introduced the Nestor algorithm
integrating Big Five personality traits and Felder-Silverman learning styles into
a Bayesian network for personalised path recommendation. Mai et al.~\cite{mai2023entropy}
demonstrated that entropy-based behavioral analysis reveals higher-performing
students exhibit lower behavioral entropy, supporting action-sequence modeling
as a diagnostic signal.

\subsection{Programming Pedagogy and Cognitive Transfer}

Scherer et al.~\cite{scherer2021transfer} documented moderate-to-large effect
sizes (Hedges' $g=0.47$--$0.73$) for cognitive transfer from programming to
mathematical reasoning. Lowe~\cite{lowe2019debugging} proposed debugging-first
pedagogy based on the Theory of Applied Mind for Programming (TAMP), arguing
that debugging engages both System~1 and System~2 thinking.

\subsection{Multi-Modal Learning Analytics and Systematic Reviews}

Wang et al.~\cite{wang2024multimodal} integrated programming activity logs and
video clickstream data into a visual analytics tool. Wang et al.~\cite{wang2024aireview}
reviewed 2,223 AI-in-education articles, identifying adaptive learning as the most
studied category while noting underutilisation of theoretical frameworks.
Zhang and Sun~\cite{zhang2024genai} identified the ``diverse demands vs.\ limited
practice'' paradox in Generative AI for education.

\subsection{Hyperspherical Representation Learning}

Han et al.~\cite{han2024hvsae} demonstrated that von Mises-Fisher distributions
over the unit sphere provide superior discrimination between similar concepts
compared to Euclidean embeddings. The CPAL system extends this to multi-modal
learner modeling, achieving $\kappa=91.1\%$ in preliminary experiments.

\subsection{Synthesis: Six Identified Gaps}

The preceding review reveals six convergent gaps: (1)~no PCM-validated LP for
code comprehension exists; (2)~informal reasoning in code explanations has never
been characterised; (3)~prompt engineering for LLM code scaffolding lacks LP
grounding; (4)~terminological ambiguity in code has never been mapped against
LLM scoring failures; (5)~the five-stage LP validity argument has never been
applied to LLM-generated outputs; (6)~equity gaps for emergent bilingual code
learners remain unstudied.

%===================================================================
\section{Methodology}
\label{sec:method}

\subsection{Architectural Overview and System Layers}

The CPAL framework follows a six-layer architecture: Presentation (FastAPI REST,
port~8000), Orchestration (Intervention Orchestrator, 11-step pipeline), Encoding
(HVSAE~+~CodeBERT~+~BERT~+~BiLSTM), Analysis (DINA~+~Nestor~+~RNN/HMM),
Reinforcement Learning (4-level Hierarchical Multi-Task RL), and Knowledge
(CSE-KG~2.0~+~Student~KG~+~Pedagogical~KG). The implementation uses Python~3.10,
PyTorch~2.0, Transformers~4.30, RDFLib~6.3, SPARQLWrapper~2.0, and
Stable-Baselines3~2.0.

Each student turn triggers the nine-component pipeline described in
Algorithm~\ref{alg:pipeline}, completing in 425--650\,ms on NVIDIA~T4 hardware.
The architecture is intentionally modular: the LLM backbone
(Groq \texttt{llama-3.1-8b-instant}), the SPARQL endpoint, and the RL policy
can each be swapped without modifying the encoding or analysis layers.

\begin{algorithm}[t]
\caption{Per-Turn Nine-Component Pipeline}
\label{alg:pipeline}
\begin{algorithmic}[1]
\REQUIRE Turn $t$: code $c_t$, error $e_t$, actions $\mathbf{a}_t$, text $x_t$
\ENSURE Response $r_t$, updated KG state
\STATE $\mathbf{z}_t \leftarrow \mathrm{HVSAE}(c_t, e_t, \mathbf{a}_t)$ \hfill [150--200\,ms]
\STATE $\boldsymbol{\alpha}_t \leftarrow \mathrm{DINA}(\mathbf{z}_t, Q)$ \hfill [50--100\,ms]
\STATE $\boldsymbol{\pi}_t \leftarrow \mathrm{Nestor}(\mathbf{a}_t, x_t)$; $\ell_t \leftarrow \mathrm{LSI}(\mathbf{a}_t, x_t)$ \hfill [75--125\,ms]
\STATE $\sigma_t \leftarrow \mathrm{COKE}(\mathbf{a}_t, \Delta\mathbf{t}_t)$ \hfill [BiLSTM 100--150\,ms]
\STATE $\mathbf{k}_t \leftarrow \mathrm{SPARQL}(\mathrm{CSE\text{-}KG}, \mathrm{concept}(\mathbf{z}_t))$ \hfill [50--75\,ms]
\STATE $m_t \leftarrow \mathrm{PedKG.detect}(c_t, e_t)$ \hfill [KG query]
\STATE $s_t \leftarrow \mathrm{CodeBERT.score}(c_t)$ \hfill [within HVSAE]
\STATE $\iota_t \leftarrow \mathrm{HierRL}(\boldsymbol{\alpha}_t, \sigma_t, \boldsymbol{\pi}_t)$
\STATE $p_t \leftarrow \mathrm{Prompt9}(\sigma_t,\boldsymbol{\pi}_t,\ell_t,m_t,\mathbf{k}_t,\iota_t,s_t)$ \hfill [100--150\,ms]
\STATE $r_t \leftarrow \mathrm{LLM}(p_t)$; $\mathrm{BidirectionalUpdate}(\mathbf{z}_t, r_t, m_t)$
\RETURN $r_t$
\end{algorithmic}
\end{algorithm}

\subsection{Datasets and Preprocessing}

Table~\ref{tab:datasets} summarises the four offline training datasets. All datasets
were accessed through their official channels and preprocessed in isolated Docker
containers to prevent data leakage between training phases.

\begin{table}[t]
\caption{Dataset Characteristics and Processing Methods}
\label{tab:datasets}
\centering
\renewcommand{\arraystretch}{1.15}
\begin{tabular}{>{\raggedright}p{1.75cm}
                >{\raggedright}p{1.7cm}
                p{4.05cm}}
\toprule
\textbf{Dataset} & \textbf{Scale} & \textbf{Processing} \\
\midrule
CodeNet & 14M submissions &
  CodeBERT tokenisation; 17-type error extraction;
  correct/buggy pairs (4--6\,h) \\[3pt]
ProgSnap2 & 4.45M events &
  Action-sequence extraction; $\Delta t$ features;
  COKE chain mining~\cite{mai2023entropy} (2--3\,h) \\[3pt]
ASSISTments & 500k+ responses &
  Q-matrix (997 skills); EM for DINA
  slip/guess~\cite{delatorre2009dina} (1--2\,h) \\[3pt]
MOOCCubeX & 1M+ interactions &
  Concept-video pairs (624k);
  prerequisite chains (3--4\,h) \\[3pt]
Evaluation & 45 turns / 10 sessions &
  Real-time; $<$500\,ms/turn;
  temporal stratification \\
\bottomrule
\end{tabular}
\end{table}

\subsubsection{CodeNet Preprocessing}
IBM CodeNet~\cite{puri2021codenet} provides 14~million code submissions across
55~languages and 4,000~problems. The pipeline extracts correct/buggy program
pairs, tokenises using the \texttt{codebert-base} vocabulary (50,265~tokens),
and classifies errors into 17~semantic types spanning logic errors, off-by-one
boundary violations, variable misuse, missing base cases, stack overflow patterns,
null/index errors, scope violations, type mismatches, and syntax misconceptions
across four severity levels.

\subsubsection{ProgSnap2 Preprocessing}
ProgSnap2~\cite{price2020progsnap2} contributes 4.45~million timestamped
programming events (EventID, SubjectID, ProblemID, EventType, CodeState,
ServerTimestamp). The pipeline groups events by (SubjectID, ProblemID) session,
extracts action sequences $\mathbf{a}=(a_1,\ldots,a_T)$, computes inter-event
time deltas $\Delta t_p$, and mines COKE cognitive-chain patterns using the
entropy-based behavioral analysis of Mai et al.~\cite{mai2023entropy}. Frustration
is operationalised as $\geq3$ consecutive \texttt{Run.Error} events with
$\sum\Delta t>60$\,s time-stuck. The derived mean time-stuck of $71.3\pm9.2$\,s
calibrates the 67.5\,s threshold in the real-time COKE module.

\subsubsection{ASSISTments Preprocessing}
The ASSISTments skill-builder dataset provides student-problem response records
(UserID, ProblemID, SkillName, Correct, HintCount, AttemptCount). Preprocessing
constructs a Q-matrix mapping 997~skill components to 4,217~problems, then
estimates DINA slip ($s_p$) and guess ($g_p$) parameters via EM following
de la Torre~\cite{delatorre2009dina}. Evidence strength weights from
\texttt{hint\_count} and \texttt{attempt\_count} penalise guess-inflated mastery.

\subsubsection{MOOCCubeX Preprocessing}
MOOCCubeX contributes 624,038~concept-video relations and 33,412~concept-problem
mappings extracted from 1M+ MOOC interaction logs. Preprocessing instantiates
prerequisite chains for the CS concept taxonomy, builds learning-progression
probability distributions per concept, and populates the Pedagogical~KG with
cross-session misconception frequency counts.

\subsection{HVSAE: Hyperspherical Variational Self-Attention Autoencoder}

The HVSAE~\cite{han2024hvsae} is the central multi-modal encoder, fusing three
parallel encoders through multi-head self-attention.

\textbf{Code Encoder.} \texttt{microsoft/codebert-base} (125M parameters)
pre-trained on CodeNet produces 768-dim code representations. Fine-tuning uses a
contrastive objective separating semantically correct from incorrect implementations.

\textbf{Error Encoder.} \texttt{bert-base-uncased} (110M parameters) fine-tuned
on ProgSnap2 error messages with cross-entropy loss over the 17-type taxonomy;
AdamW, $\text{lr}=2\times10^{-5}$, $\epsilon=10^{-8}$, 3~epochs.

\textbf{Behavioral Encoder.} A bidirectional LSTM (2~layers, 256~hidden units per
direction, dropout~0.2) with additive attention trained on action tuples
$(a_p, \Delta t_p, o_p)$ from ProgSnap2. Higher attention weights are assigned to
time-stuck events, operationalising the frustration signal of Mai et al.~\cite{mai2023entropy}.

Modality fusion uses 8-head self-attention with $d_k=64$:
\begin{equation}
\text{Attention}(Q,K,V)=\text{softmax}\!\left(\tfrac{QK^\top}{\sqrt{d_k}}\right)V
\label{eq:attention}
\end{equation}

The fused 512-dim vector is projected to the unit hypersphere via the von
Mises-Fisher (vMF) distribution:
\begin{equation}
f(\mathbf{x}\mid\boldsymbol{\mu},\kappa)=C_d(\kappa)\exp(\kappa\boldsymbol{\mu}^\top\mathbf{x})
\label{eq:vmf}
\end{equation}
where $C_d(\kappa)=\kappa^{d/2-1}\big/\!\left((2\pi)^{d/2}I_{d/2-1}(\kappa)\right)$.
Higher $\kappa$ indicates more discriminative representations.

The composite training objective is:
\begin{equation}
\mathcal{L}=\lambda_1\mathcal{L}_{\text{recon}}+\lambda_2\mathcal{L}_{\text{KL}}
+\lambda_3\mathcal{L}_{\text{cont}},
\;\;(\lambda_1,\lambda_2,\lambda_3)=(0.5,0.3,0.2)
\label{eq:loss}
\end{equation}
Training used AdamW ($\text{lr}=2\times10^{-4}$, weight decay~0.01, cosine
annealing, 30~epochs, batch size~64, 10\% warm-up). An MLP misconception classifier
on the 256-dim bottleneck outputs softmax probabilities over five types:
(1)~missing base case, (2)~off-by-one errors, (3)~variable misuse,
(4)~algorithm misunderstanding, (5)~syntax misconception.

\subsection{COKE Cognitive-State Tracking}

COKE infers student cognitive state from action tuples $(a_p,\Delta t_p,o_p)$ where
$a_p$ is action type, $\Delta t_p$ is elapsed time, and
$o_p\in\{\text{success,error,neutral}\}$. Four states are tracked:
\textit{perceiving}, \textit{confused}, \textit{frustrated}, \textit{understanding},
stored in a \texttt{CognitiveChain} directed graph with per-turn confidence and
frequency updates.

The frustration heuristic fires when the sequence contains $\geq3$ consecutive
\texttt{run\_test} events with cumulative time-stuck $\geq67.5$\,s and no success
outcome. The \texttt{frustrated} state triggers: (1)~de-escalating tone selection;
(2)~one-level reduction in scaffolding complexity; (3)~emotional weight increase
in the RL objective to $\geq50\%$ (Section~\ref{sec:rl}).

In the traced session (\texttt{student\_sample\_03}), the five-turn COKE trajectory
was: \textit{frustrated}~$\to$~\textit{confused}~$\to$~\textit{engaged}~$\to$
\textit{understanding}~$\to$~\textit{confused}. The terminal dip corresponds to
introduction of memoisation as a genuinely new concept, confirming correct scaffolding
behaviour: the system did not assume recursion mastery transfers automatically.

\subsection{Nestor Bayesian Personality Profiling}

The Nestor module~\cite{nadimpalli2025nestor} infers Big Five traits (OCEAN) and
Felder-Silverman learning styles via:
\begin{equation}
P(\text{Intervention}\mid\mathbf{E})\approx\text{softmax}
(W_2\cdot\text{ReLU}(W_1\cdot\mathbf{E}+b_1)+b_2)
\label{eq:nestor}
\end{equation}
Parameters minimise KL divergence from exact Bayesian inference on 10,000~synthetic
student sequences.

Dynamic learning-style inference combines three evidence sources in priority order:
(1)~\textit{stored cross-session profile} from the Student~KG;
(2)~\textit{behavioral inference}---visual vs.\ verbal inferred from debugger
vs.\ print-statement frequency; active vs.\ reflective from time-before-first-run
($<$30\,s = active; $\geq$30\,s = reflective); sequential vs.\ global from
edit-run pair density;
(3)~\textit{NLP keyword analysis}---visual (\textit{diagram, picture, show}),
active (\textit{try, practice, run}), sequential (\textit{step, first, then}).

At Turn~1 of the traced session, Nestor inferred Openness~$=0.58$,
Extraversion~$=0.34$, Neuroticism~$=0.57$, prescribing: ``give examples first;
do not assume follow-up questions; use supportive tone.'' This drove a response
opening with ``Don't worry'' and an ASCII call-stack recursion diagram (visual style
confirmed from chat keywords).

\subsection{DINA Cognitive Diagnosis}

The DINA model~\cite{delatorre2009dina} provides interpretable knowledge-component
mastery estimation:
\begin{equation}
P(Y_{sp}=1\mid\boldsymbol{\alpha}_s)=(1-s_p)^{\eta_{sp}}g_p^{1-\eta_{sp}},
\quad\eta_{sp}=\prod_{k=1}^{K}\alpha_{sk}^{q_{pk}}
\label{eq:dina}
\end{equation}
where $\alpha_{sk}\in\{0,1\}$ is the latent mastery indicator, $q_{pk}$ is the
Q-matrix entry, $s_p$ is slip probability, and $g_p$ is guess probability.
Components with mastery $\hat{\alpha}_{sk}<0.70$ surface as \textit{knowledge\_gap}
nodes in the Student~KG and trigger CSE-KG~2.0 prerequisite-path queries included
in the nine-section prompt.

\subsection{Hierarchical Multi-Task Reinforcement Learning}
\label{sec:rl}

A four-level Hierarchical Multi-Task RL controller (PPO via Stable-Baselines3~2.0)
governs intervention selection~\cite{wan2023rl}:

\textbf{Level~1 (Meta-Controller):} Learns general teaching strategy for the student
type over the full curriculum horizon.

\textbf{Level~2 (Curriculum Controller):} Selects which concept to address next,
balancing prerequisite satisfaction against student engagement.

\textbf{Level~3 (Session Multi-Task Controller):} Simultaneously optimises five
objectives with adaptive weights: learning maximisation (40\% baseline), engagement
optimisation (25\%), emotional state management (20\%), time efficiency (10\%),
and long-term retention (5\%). When $\sigma_t=\text{frustrated}$, emotional weight
increases to $\geq50\%$; when $\sigma_t=\text{engaged}$, learning weight increases
to $\geq60\%$.

\textbf{Level~4 (Intervention Executor):} Maps the Level~3 decision to one of five
intervention types: \textit{visual\_explanation}, \textit{interactive\_exercise},
\textit{guided\_practice}, \textit{conceptual\_deepdive},
\textit{motivational\_support}.

The reward function combines immediate mastery gain $\Delta\hat{\alpha}$, emotion
transition quality (frustration-reduction bonus), engagement score, and a retention
proxy (spaced-repetition alignment). Policies use PPO with $\gamma=0.99$,
GAE $\lambda=0.95$, and a replay buffer of 10,000~transitions.

\subsection{Knowledge Graph Integration}

Three knowledge graphs are maintained and bidirectionally
updated~\cite{abusalih2024kg,shengju2025cross}.

\textbf{CSE-KG~2.0}~\cite{dessi2025csekg}: A live RDF/OWL domain knowledge base
queried via SPARQL. Entity types include \texttt{cskg:Concept}, \texttt{cskg:Method},
and \texttt{cskg:Task} with \texttt{requiresKnowledge} prerequisite edges.

\textbf{Student KG}: A per-student directed graph tracking mastery-state nodes
($\hat{\alpha}_{sk}$), misconception-pattern nodes, and cross-session
learning-trajectory edges, updated bottom-up from behavioral data after each turn.

\textbf{Pedagogical KG}: A misconception ontology with severity and frequency
counters updated cross-session. Each node carries: concept link, severity
(\textit{critical/high/medium}), detection frequency, correction strategy, and
recommended intervention type. Across 10~evaluation conversations, the KG
accumulated four persistent nodes including \texttt{mc\_recursion\_no\_base\_case}
(frequency~$=0.86$, severity~$=\textit{critical}$) and \texttt{mc\_arrays\_indexerror}
(frequency updated $0.28\!\to\!0.29$ on second occurrence).

Updates follow a bidirectional mechanism: bottom-up (behavioral data $\to$
Student~KG) and top-down (curriculum objectives $\to$ intervention selection).

\subsection{Nine-Section Prompt and Ten Personalization Features}
\label{sec:prompt}

Each turn produces a 2,824--3,043-character nine-section structured prompt
implementing the Model~B context strategy of Jin et al.~\cite{jin2025automated}.
The nine sections encode: (1)~COKE state with confidence; (2)~Nestor Big Five
vector; (3)~Felder-Silverman learning style with behavioral evidence; (4)~detected
misconception type and correction strategy; (5)~CSE-KG prerequisite gap chain
(mastery~$<0.70$); (6)~RL-selected intervention type; (7)~prior three-turn context;
(8)~CodeBERT correctness score; (9)~LP-level target and rubric for the current topic.

These sections drive ten adaptive personalization features at generation time:
(F1)~Conversation Memory; (F2)~Emotional Tone Adaptation (gentle/supportive/
enthusiastic/celebratory driven by frustration and engagement thresholds);
(F3)~Learning Style Personalisation (visual analogies, numbered steps, hands-on
exercises, or conceptual explanations matched to Nestor output);
(F4)~Personality-Based Communication (extra reassurance for high Neuroticism;
structured format for high Conscientiousness; advanced connections for high
Openness); (F5)~Progress-Aware Scaffolding; (F6)~Interest-Based Examples;
(F7)~Format Preference Adaptation (ASCII diagrams, code blocks, bulleted steps);
(F8)~Error-Specific Feedback from Pedagogical~KG;
(F9)~Metacognitive Guidance calibrated to LP level; and
(F10)~Difficulty Adaptation via mastery gap and COKE state.

\subsection{Three-Tier Evaluation Framework}

\textbf{Tier~1 (Keyword Baseline):} Concept detection via keyword matching with
morphological variation and WordNet synonym expansion; error detection maps
17~error types to keyword patterns; six binary quality indicators.

\textbf{Tier~2 (HVSAE-Enhanced Semantic):} Cosine distance in the 256-dim
hyperspherical space; MLP misconception classifier; quality boost~$=\kappa\times0.3$.

\textbf{Tier~3 (Standard NLP):} BLEU, ROUGE-1/2/L, and METEOR for comparison
with other educational dialogue systems.

The overall score integrates all three tiers:
\begin{equation}
\text{Score}=0.4C+0.3E+0.2Q+0.1\overline{\text{NLP}}
\label{eq:score}
\end{equation}
where $C=$ concept coverage, $E=$ error coverage, $Q=$ quality score with HVSAE
boost, $\overline{\text{NLP}}=$ mean of BLEU, ROUGE-L, and METEOR.

\subsection{Training Hyperparameters and Latency Profile}

\begin{table}[t]
\caption{Component Latency Per Student Turn (NVIDIA T4, 16\,GB)}
\label{tab:latency}
\centering
\begin{tabular}{lc}
\toprule
\textbf{Component} & \textbf{Latency (ms)} \\
\midrule
HVSAE multi-modal encoding & 150--200 \\
DINA cognitive diagnosis & 50--100 \\
Nestor psychological profiling & 75--125 \\
Behavioral RNN / HMM & 100--150 \\
KG query and update & 50--75 \\
Prompt build + LLM call & 100--150 \\
\midrule
\textbf{Total per turn} & \textbf{425--650} \\
\bottomrule
\end{tabular}
\end{table}

HVSAE: AdamW, $\text{lr}=2\times10^{-4}$, weight decay~0.01, $\beta_1=0.9$,
$\beta_2=0.999$, batch~64, 30~epochs, cosine annealing with 10\% warm-up.
Behavioral LSTM: SGD with momentum~0.9, cyclic LR (base $10^{-4}$, max $10^{-2}$).
DINA~EM: converged in $\leq200$ iterations for all 997~skill components.
RL: PPO with $\gamma=0.99$, GAE $\lambda=0.95$, replay buffer of 10,000~transitions.
Table~\ref{tab:latency} reports latency profiling on the evaluation hardware.

%===================================================================
\section{Results and Findings}
\label{sec:results}

\subsection{Aggregate Performance}

Table~\ref{tab:agg} reports performance across 10~conversations (45~turns). The
42.8\% accuracy improvement from Tier~1 to Tier~2 directly mirrors the Model~A
to Model~B QWK gains of Jin et al.~\cite{jin2025automated}, providing preliminary
support for RQ3. Misconception detection confidence of 52.7\% (SD~$=12.4\%$)
indicates moderate uncertainty: since $\kappa=91.1\%$ confirms the latent space
is highly discriminative, the source of uncertainty lies at decision boundaries
between misconception types sharing informal linguistic surface forms---consistent
with the RQ2 hypothesis.

\begin{table}[t]
\caption{Aggregate System Performance ($N=10$, 45~Turns)}
\label{tab:agg}
\centering
\renewcommand{\arraystretch}{1.1}
\begin{tabular}{lcc}
\toprule
\textbf{Metric} & \textbf{Value} & \textbf{SD} \\
\midrule
Overall accuracy (Tier~2) & 49.7\% & 4.2\% \\
Tier~1 baseline accuracy & 34.8\% & 5.1\% \\
HVSAE improvement & $+$42.8\% & --- \\
Error accuracy & 100.0\% & 0.0\% \\
Error precision / recall & 100.0\% & 0.0\% \\
Concept detection rate & 62.3\% & 8.7\% \\
Quality score & 48.7\% & 25.3\% \\
Emotion matching accuracy & 100.0\% & 0.0\% \\
Encoding quality ($\kappa$) & 91.1\% & 3.2\% \\
Misconception confidence & 52.7\% & 12.4\% \\
Semantic similarity & 0.78 & 0.09 \\
\bottomrule
\end{tabular}
\end{table}

\subsection{Quality Indicator and Misconception Analysis}

Only 48.9\% of responses included code examples and 40.0\% used structural
elements (Table~\ref{tab:quality})---both strong indicators of pedagogical quality
in programming education~\cite{scherer2021transfer}. Improving code-example
generation from 48.9\% to $\geq70\%$ is the primary engineering target.

\begin{table}[t]
\caption{Quality Indicators and Misconception Detection Gains}
\label{tab:quality}
\centering
\renewcommand{\arraystretch}{1.1}
\begin{tabular}{lcc}
\toprule
\textbf{Indicator / Misconception Type} & \textbf{Baseline} & \textbf{HVSAE} \\
\midrule
Code example present & 48.9\% & --- \\
Explanation (why, not just what) & 60.0\% & --- \\
Actionable solution & 71.1\% & --- \\
Structure (headings / bullets) & 40.0\% & --- \\
\midrule
Missing base case & 82\% & 91\% \\
Off-by-one errors & 65\% & 73\% \\
Variable misuse & 58\% & 66\% \\
Algorithm misunderstanding & 47\% & 55\% \\
Syntax misconception & 41\% & 49\% \\
\bottomrule
\end{tabular}
\end{table}

Greater relative gains on conceptual errors (base case, off-by-one) than syntactic
errors confirm that hyperspherical encoding adds most value for semantically complex
misconceptions, providing empirical direction for RQ2.

\subsection{Single-Student Trace and RL Validation}

The five-turn code-correctness trajectory for \texttt{student\_sample\_03} was:
0.70 (T1: infinite recursion) $\to$ 0.84 (T2: base case added) $\to$ 0.96
(T3: Fibonacci) $\to$ 0.84 (T4--T5: memoisation). The terminal correctness dip
and COKE \textit{understanding}~$\to$~\textit{confused} transition on Turn~5
confirm correct scaffolding: the system did not assume recursion mastery transfers
to memoisation. The Level~3 RL controller correctly increased emotional weight
from 20\% to 35\% between Turns~1--2 (frustration), then reduced it to 15\% by
Turn~3 as COKE transitioned to \textit{engaged}, demonstrating adaptive multi-task
weight adjustment.

\subsection{Performance Distribution and Comparative Positioning}

\begin{table}[t]
\caption{Performance Distribution (45 Evaluation Turns)}
\label{tab:dist}
\centering
\begin{tabular}{lccc}
\toprule
\textbf{Tier} & \textbf{Turns} & \textbf{\%} & \textbf{Interpretation} \\
\midrule
Excellent ($\geq$80\%) & 0 & 0.0\% & Not yet achieved \\
Good (60--80\%) & 8 & 17.8\% & Strong error detection \\
Fair (40--60\%) & 37 & 82.2\% & Consistent baseline \\
Poor ($<$40\%) & 0 & 0.0\% & Robust lower bound \\
\bottomrule
\end{tabular}
\end{table}

\begin{table}[t]
\caption{Comparative Analysis with Representative Systems}
\label{tab:comp}
\centering
\renewcommand{\arraystretch}{1.1}
\begin{tabular}{>{\raggedright}p{2.0cm}ccp{1.8cm}}
\toprule
\textbf{System} & \textbf{Err.} & \textbf{Pers.} & \textbf{Limitation} \\
\midrule
CPAL (Ours) & 100\% & Multi-dim. & Concept artifact \\
Nestor BN~\cite{nadimpalli2025nestor} & N/A & Psych. only & Survey dep. \\
KG Approaches~\cite{abusalih2024kg} & 85--90\% & KG-based & Static KG \\
RL Path~\cite{wan2023rl} & N/A & Sequential & No error analysis \\
Traditional ITS & 70--85\% & Limited & Single-modality \\
\bottomrule
\end{tabular}
\end{table}

Concentration in the Fair tier (82.2\%) with zero Poor turns confirms a robust
baseline. The concept coverage data artifact (0.0\% evaluated accuracy alongside
62.3\% system detection rate) indicates missing expected-concept annotation tags
in the evaluation harness, not a pipeline failure; resolving this is expected to
raise Tier~2 accuracy from 49.7\% to the 60--70\% range.

%===================================================================
\section{Remaining Work and Timeline}
\label{sec:remaining}

The remaining work implements and extends the LLM Explanation Evaluation
Methodology~\cite{patel2025llmeval}.

\subsection{Phase~1: PCM-Validated Learning Progression (Weeks~1--4)}

\textbf{Milestone~1.1} (Weeks~1--2): Think-aloud interviews; 60+ written constructed
responses from 15+ students across three experience levels on recursion, iteration,
and scope.

\textbf{Milestone~1.2} (Weeks~2--3): Grounded theory coding deriving four LP
levels: (L1)~Heuristic---``it calls itself until it stops''; (L2)~Qualitative---
``each call breaks the problem into a smaller version''; (L3)~Measurable Variables---
``base case stops recursion; each frame holds local state''; (L4)~Scientific
Explanation---$T(n)=T(n{-}1)+O(1)$; space $O(n)$.

\textbf{Milestone~1.3} (Weeks~3--4): PCM threshold ordering via R \texttt{eRm}
following Jin et al.~\cite{jin2019validation}. Disordered thresholds require LP
revision before any subsequent milestone (non-negotiable gate for RQ1).

\subsection{Phase~2: LLM Comparison Experiment (Weeks~4--9)}

Four LLMs are compared under Model~A (question only) and Model~B (full nine-section
CPAL context): GPT-4o, Claude~3.5 Sonnet, Gemini~1.5 Pro, and Llama-3.1-8b.
The evaluation engine is the Anthropic API (\texttt{claude-sonnet-4-20250514}).

The composite analytic score~\cite{patel2025llmeval} is:
\begin{equation}
S_{\text{analytic}}=0.35P+0.25L+0.25A+0.15C
\label{eq:analytic}
\end{equation}
where $P=$ Pedagogical Value, $L=$ Language Clarity, $A=$ Technical Accuracy,
$C=$ Context Alignment. Explanations are ranked by: LP Level $\succ$
$S_{\text{analytic}}$ $\succ$ Error Flag Count. QWK is computed relative to an
ideal Level~4 expert explanation using Williamson et al.\ thresholds~\cite{williamson2012framework}:
$\geq0.85$ excellent, $0.70$--$0.85$ strong, $0.60$--$0.70$ moderate, $<0.60$ poor.
Five error flags---Ambiguous Terminology, Logical Inconsistency, Jargon Without
Grounding, Missing Contextual Grounding, Informal Language Obscuring Meaning---can
each reduce the effective LP level.

\subsection{Phase~3: Equity Analysis and Validity Argument (Weeks~9--12)}

Informal language corpus construction (RQ2); terminological ambiguity QWK mapping
on code synonym pairs (``higher-order function'' vs.\ ``function taking a function'';
``closure'' vs.\ ``lexical scope''; ``memoisation'' vs.\ ``caching''); DIF analysis
across bilingual/experience subgroups (RQ4); five-stage validity argument for the
winning LLM; paper writing targeting EDM~2026, AIED~2026, SIGCSE~2026.

\begin{table}[t]
\caption{Remaining Work Timeline (Twelve Weeks)}
\label{tab:timeline}
\centering
\begin{tabular}{cp{6.3cm}}
\toprule
\textbf{Week(s)} & \textbf{Milestone} \\
\midrule
1--2 & Clinical interviews; 60+ constructed responses \\
2--3 & Thematic analysis; four-level code-LP derived \\
3--4 & PCM threshold ordering; LP revision if needed (RQ1) \\
4--5 & Model A/B conditions; explanations from 4~LLMs \\
5--6 & LP level assigned per explanation (5 decision rules) \\
6--7 & Analytic scoring ($P/L/A/C$); five error flags applied \\
7--9 & QWK; three-tier ranking; winner identified (RQ3) \\
9--10 & Informal language corpus; misclassification rate (RQ2) \\
10 & Terminological ambiguity QWK degradation \\
10--11 & DIF analysis across subgroups (RQ4) \\
11 & Five-stage validity argument for winning LLM \\
12 & Paper writing; advisor review; submission \\
\bottomrule
\end{tabular}
\end{table}

%===================================================================
\section{Conclusion}
\label{sec:conclusion}

This midterm report presents a framework bridging Learning Progression theory,
automated LP-based scoring, and LLM evaluation methodology. The CPAL
implementation---a nine-component, ten-feature, six-layer system with four-level
hierarchical RL and three dynamically updated knowledge graphs---is fully
operational. Preliminary results across 45~interaction turns demonstrate perfect
error and emotion detection (100\%), a 42.8\% accuracy improvement from LP-grounded
context provision, and highly discriminative hyperspherical representations
($\kappa=91.1\%$). Four Research Questions structure a twelve-week remaining
programme culminating in PCM validation, a five-LLM comparison experiment, DIF
equity analysis, and a five-stage validity argument for LLM-generated code
explanations. The goal is AI tools that can be \textit{proven} to develop
students' understanding of code---and the psychometric evidence to substantiate
that claim.

%===================================================================
\bibliographystyle{IEEEtran}
\bibliography{references}

\end{document}
